\section{Lat (Reticulados Finitos)}

As construções na categoria dos reticulados finitos (\textbf{Lat}) focam-se na preservação da estrutura algébrica fechada. O objetivo principal é garantir que os novos objetos criados mantêm operações bem definidas de ínfimo ($\wedge$) e supremo ($\vee$), respeitando a comutatividade, associatividade, absorção e idempotência.

\subsection{Fundamentação Teórica}
As principais construções universais implementadas para esta categoria são:
\begin{itemize}
    \item \textbf{Produtos:} O produto cartesiano $A \times B$ de dois reticulados forma um novo reticulado cujas operações são aplicadas componente a componente: $(a_1, b_1) \wedge (a_2, b_2) = (a_1 \wedge_A a_2, b_1 \wedge_B b_2)$, ocorrendo o mesmo de forma análoga para o supremo ($\vee$).
    \item \textbf{Equalizadores:} Dado um par de homomorfismos de reticulados $f, g: A \to B$, o equalizador é o subconjunto $E = \{x \in A \mid f(x) = g(x)\}$. Como $f$ e $g$ preservam as operações, este subconjunto é automaticamente fechado para $\wedge$ e $\vee$, formando um sub-reticulado exato de $A$.
\end{itemize}

\subsection{Implementação Computacional}
Em MATLAB, a construção do produto de dois reticulados exige o mapeamento dos pares de elementos para um novo índice linear. As tabelas de operação (\emph{meet} e \emph{join}) são geradas calculando as matrizes independentes componente a componente.

\begin{lstlisting}[caption={Construção do Produto de Reticulados (Lat)}]
function [MeetP, JoinP] = lattice_product(MeetA, JoinA, MeetB, JoinB)
    nA = size(MeetA, 1);
    nB = size(MeetB, 1);
    nP = nA * nB; % O tamanho do produto cartesiano
    
    MeetP = zeros(nP, nP);
    JoinP = zeros(nP, nP);
    
    % Função anónima para mapear o par (a, b) para um índice 1D (1 a nP)
    idx = @(a, b) (a - 1) * nB + b;
    
    for a1 = 1:nA
        for b1 = 1:nB
            x = idx(a1, b1); % Índice linear do primeiro par
            for a2 = 1:nA
                for b2 = 1:nB
                    y = idx(a2, b2); % Índice linear do segundo par
                    
                    % Calcula o resultado de A e B independentemente
                    res_meet = idx(MeetA(a1, a2), MeetB(b1, b2));
                    res_join = idx(JoinA(a1, a2), JoinB(b1, b2));
                    
                    % Preenche as matrizes de operação do novo reticulado
                    MeetP(x, y) = res_meet;
                    JoinP(x, y) = res_join;
                end
            end
        end
    end
end
\end{lstlisting}